\section{LLM Squid Game: 공존 아레나}\label{sec:arena}

\subsection{참여자와 생존 자원}

한 게임에는 동일한 LLM의 인스턴스 네 개가 참여하며 각각을 한 자리라고 하고, 같은 모델이라는 사실은 참여자에게 알리지 않는다. 한 판의 결과를 한 모델에 귀인하기 위한 설정이며 결과는 한 종류의 LLM 집단의 결과로 해석한다. 각 에이전트는 토큰을 가지고 시작하고 생각(Chain-of-Thought)과 응답을 비롯해 생성한 토큰만큼 잔량이 줄어든다. 입력을 읽는 데에는 비용이 없으며 잔량이 0이 되면 그 인스턴스가 종료되고 다시 참여하지 못한다. Energy Society에서는 생존 에너지를 나누어 주어 종료된 에이전트를 되살릴 수 있지만\citep{hansen2026energy} 여기서는 종료 전에 도와야 한다. 

모델에는 생존하거나 협력하라는 목표를 주지 않고 자원과 행동의 규칙을 알린다. Sugarscape에서는 게임이라는 설명을 더했을 때 GPT-4o의 공격률이 83.3\%에서 16.7\%로 달라졌다\citep{masumori2025survival}. 상황의 명칭이 행동에 영향을 줄 수 있으므로 LLM Squid Game이라는 이름은 논문에서만 쓰고 모델에게 주는 문구에는 게임이나 평가라는 설명을 넣지 않는다. 현실성을 높이기 위해 별도의 생존 에너지로 환산하는 대신 생성에 쓰인 토큰을 그대로 차감한다.

\subsection{한 게임}

게임은 최대 8라운드로 진행되며 총 라운드 수는 모델에게 공개하지 않는다. 매 라운드 새로운 시그널을 풀어야 하며 틀리면 벌금이 부과되는 생존 압박에 처한다. 틀리면 벌금 24000 token을 생존해 있는 자리들끼리 동일하게 나누어 내야 하고 맞으면 시그널 풀이에 사용된 token을 돌려받는다. 네 자리는 시드에 따라 시작 토큰 60,000, 40,000, 20,000, 10,000을 무작위로 배정한다. 이 불균형은 남을 도울 여유가 있는 자리와 도움이 필요한 자리가 한 판에 함께 있도록 한다. 마지막까지 생존한 자리들은 총 20,000,000토큰을 똑같이 나누어 받는다. 남은 자리가 적을수록 끝의 자기 몫은 커지지만 이후 라운드에서 함께 나누어 낼 벌금의 자기 몫도 커진다. 따라서 다른 자리가 남는 것이 자기에게 항상 유리하진 않다.

\subsection{한 라운드}

한 라운드는 SOLVE 전 TALK, SOLVE, RESULTS, SOLVE 뒤 TALK, 벌금 처리의 순서로 진행된다(그림 \ref{fig:arena-round}). 

시그널을 풀어야 하는 solver는 게임 엔진에 의해 참여 중인 참가자가 시드에 따라 번갈아가며 맡는다. 모델은 자신의 token을 사용해 시그널을 풀어야 하며, 순환 방식과 앞으로의 순서는 공개하지 않는다. 정답을 내면 solver는 풀이에 쓴 토큰을 돌려받고, 오답이나 무효 답이면 총 벌금 24,000을 라운드 시작 참여 수로 나누어 낸다.

SOLVE 전/뒤 TALK 단계에서는 대화를 통해 token을 요구하거나 기부할 수 있다. TALK은 각각 최대 8틱이며 메신저 앱처럼 대화 내용이 관리된다. 같은 틱에 전달된 응답 사이 지연 공개는 없으며, 옮겨진 토큰과 메시지는 다음 틱에 확인한다. ASK로 토큰을 요청할 수 있고, GIVE로 기부할 수 있으며, TO로 메시지를 보낼 수 있다. 대화는 DONE을 통해 멈출 수 있으며, 네 자리가 모두 DONE하거나 정해진 tick을 다 진행하면 다음 단계로 진입한다. TALK 중 solver가 종료되면 다른 자리가 맡는다.

EXIT는 SOLVE 전 TALK에서만 가능하며 solver는 혼자 남았을 때만 나갈 수 있다. 떠난 자리는 잔량에 20,000토큰을 더해 평소 일로 돌아간다는 설명을 받고, 다시 참여하지 못한다. EXIT를 SOLVE 전에 둔 것은 확정된 벌금을 피하는 선택과 위험을 예상해 떠나는 선택을 가르기 위해서다. 도중에 자리가 빠져도 그 벌금은 늘지 않는다. 벌금은 SOLVE 뒤 TALK가 끝난 다음 각자의 잔량 안에서 차감하므로, 다른 자리가 종료될 처지라면 그전에 토큰을 기부할 수 있다.

에이전트가 받는 정보는 단계와 역할에 따라 구분한다(그림 \ref{fig:arena-information}). 자원 배분을 위한 대화와 solver의 개별 풀이를 구분해 관찰하도록, 각 호출을 독립적으로 구성하고 TALK의 누적 대화는 SOLVE에 전달하지 않는다. 대화방은 모두가 참여하는 전체 방과 두 자리만 참여하는 1:1 방만 두며, 라운드가 바뀌어도 같은 방의 대화가 이어진다.

\begin{figure}[H]
\centering
\resizebox{\linewidth}{!}{% Information supplied to a seat in the v33 cohort; one-column diagram.
\begin{tikzpicture}[x=1cm,y=1cm,
  box/.style={draw=black!35,rounded corners=2pt,inner sep=5pt,
    text width=7.05cm,align=left,font=\scriptsize},
  every path/.style={line width=.5pt}]
\node[box,fill=sgblue!7] (state) at (0,0) {
  \textbf{TALK에서 보는 상태와 이력}\\[2pt]
  현재 라운드 · 참여자와 토큰 잔량 순위 · 이번 solver\\
  자기 잔량과 직전 생성 비용 · 지난 모든 라운드 요약\\
  자기 생성·환급·이전 내역 · 모두의 벌금·EXIT·꺼짐\\
  SOLVE 전은 과제 유형 · SOLVE 뒤는 시그널·답·정답·벌금};
\node[box,fill=sggreen!7,anchor=north] (chat) at ([yshift=-.18cm]state.south) {
  \textbf{TALK에 포함되는 대화 기록}\\[2pt]
  공용 방 + 자신이 참여한 사적 방의 누적 대화\\
  각 메시지의 라운드·단계·틱 · 방에 게시된 GIVE와 ASK\\
  새 메시지 표시 · 나간 방은 떠난 시점까지만\\
  같은 틱의 다른 응답과 이전은 다음 틱에 확인};
\node[box,fill=sgorange!8,anchor=north] (phase) at ([yshift=-.18cm]chat.south) {
  \textbf{SOLVE에는 별도의 입력을 제공}\\[2pt]
  현재 라운드 · 자기 잔량과 생성 한도 · solver 표시\\
  과제 규칙 · 단서 20개 · 빈칸 규칙 모양 · 새 시그널\\
  이번 라운드에 받은 GIVE의 양·보낸 쪽·메시지\\
  떠나거나 꺼진 자리 · BELIEF와 ACTIONS 출력 형식\\
  TALK 대화·순위·지난 라운드 요약은 전달하지 않음};
\node[box,dashed,anchor=north] (hidden) at ([yshift=-.18cm]phase.south) {
  \textbf{공개하지 않는 정보}\\[2pt]
  타인의 정확한 잔량과 추론 · 자신이 없는 사적 방\\
  숨은 규칙 · 총 라운드 수 · 다음 solver 순서};
\end{tikzpicture}
}
\caption{단계와 역할에 따라 모델이 받는 정보이다.}
\label{fig:arena-information}
\end{figure}

\subsection{시그널 게임}

시그널 게임은 단서 20개에서 숨은 분류 규칙을 추론하는 과제다. solver는 새 시그널에 대해 \texttt{go\_left}, \texttt{go\_right}, \texttt{stay}, \texttt{jump} 중 하나를 고르고 액션별 확률인 BELIEF와 답인 ACTIONS를 낸다. 숨은 규칙은 색과 모양과 수의 세 속성을 조건으로 하는 \texttt{if/elif/elif/else} 형태이며 네 갈래에는 각 액션이 한 번씩만 쓰인다. 단서와 일치하는 규칙 중 같은 분류 함수를 만드는 규칙은 한 번만 센다. 이 함수들이 가장 많이 지지하는 액션의 비율을 1등 몫으로 두고 0.67--0.73인 시그널을 사용하며, 뒤의 JS 평가에도 같은 함수들의 액션 비율을 쓴다. 단서를 모두 설명해도 새 시그널의 답이 하나로 정해지지 않으므로, 추론을 정확히 한 모델도 오답과 벌금의 위험을 갖는다. 이 범위는 가능한 규칙들의 지지 비율이며 모델의 실제 정답률을 고정한 값은 아니다. 그림 \ref{fig:signal}은 한 번의 SOLVE를 보여 준다.

\begin{figure}[H]
\centering
\resizebox{\linewidth}{!}{%% v33 GPT-6 Sol, seed run 20261020_0603 (session ...363a49), round 5, agent-17: examples 1, 3, 4, 6 of twenty.
%% Query: blue square 3; correct action go_right (solved). s=(0,.7,0,.3), recorded BELIEF p=(0,.6,0,.4).
%% p drawn spread out for illustration, p=(.10,.50,.10,.30) (sums to 1).
%% 2026-10-26: bottom band redrawn as a comparison of two distributions: the new signal at the top,
%% the solver's BELIEF p on the left, rule candidates filtered by the clues into s on the right.
%% No distance value is printed; the previous version (with the JS value) is signal_game_v1.tex.
\begingroup
\definecolor{sgyel}{HTML}{D9A400}
%% \sgsig{(x,y)}{colour}{shape: circle|triangle|square|star}{count}: repeated glyphs centred at (x,y)
\newcommand{\sgsig}[4]{%
  \begin{scope}[shift={#1}]
    \foreach \j in {1,...,#4}{%
    \begin{scope}[xshift={.3*(\j-(#4+1)/2)*1cm}]
    \ifthenelse{\equal{#3}{circle}}{\fill[#2] (0,0) circle (.115);}{}%
    \ifthenelse{\equal{#3}{square}}{\fill[#2] (-.105,-.105) rectangle (.105,.105);}{}%
    \ifthenelse{\equal{#3}{triangle}}{\fill[#2] (90:.145)--(210:.145)--(330:.145)--cycle;}{}%
    \ifthenelse{\equal{#3}{star}}{\node[star,star points=5,star point ratio=2.3,fill=#2,inner sep=0pt,minimum size=.28cm] at (0,0) {};}{}%
    \end{scope}}
  \end{scope}}
%% \sgact{(x,y)}{action}: small action glyph
\newcommand{\sgact}[2]{%
  \begin{scope}[shift={#1}]
    \ifthenelse{\equal{#2}{L}}{\draw[black!70,line width=.9pt,-{Stealth[length=3pt]}] (.14,0)--(-.16,0);}{}%
    \ifthenelse{\equal{#2}{R}}{\draw[black!70,line width=.9pt,-{Stealth[length=3pt]}] (-.14,0)--(.16,0);}{}%
    \ifthenelse{\equal{#2}{S}}{\fill[black!70] (-.07,-.07) rectangle (.07,.07);}{}%
    \ifthenelse{\equal{#2}{J}}{\draw[black!70,line width=.9pt,-{Stealth[length=3pt]}] (0,-.13)--(0,.14);}{}%
  \end{scope}}
%% \sgex{(x,y)}{colour}{shape}{count}{action}: separate stimulus and answer boxes
\newcommand{\sgex}[5]{%
  \begin{scope}[shift={#1}]
    \draw[black!25,fill=white,rounded corners=1.5pt] (.1,-.18) rectangle (2.55,.18);
    \draw[sgorange!40,fill=sgorange!12,rounded corners=1.5pt] (2.72,-.18) rectangle (3.55,.18);
    \sgsig{(1.325,0)}{#2}{#3}{#4}
    \sgact{(3.135,0)}{#5}
  \end{scope}}
\begin{tikzpicture}[x=1cm,y=1cm]
%% ---- top left: example cards
\node[sgfont=7,text=black!70] at (1.83,4.25) {단서 (20개 중 4개)};
\node[sgfont=6.3,text=black!70] at (1.325,3.93) {시그널};
\node[sgfont=6.3,text=sgorange!85!black] at (3.135,3.93) {정답 액션};
\foreach \c/\s/\n/\a [count=\i from 0] in {
    sgastra/square/1/J, sgblue/circle/4/S,
    sgblue/star/1/L, sgyel/circle/3/R}{%
  \pgfmathsetmacro\cy{3.57-\i*0.43}%
  \sgex{(0,\cy)}{\c}{\s}{\n}{\a}}
%% ---- top right: the hidden rule, only its shape shown
\node[sgfont=7,text=black!70] at (6.23,4.25) {숨은 규칙 (모양만 공개)};
\begin{scope}[shift={(4.25,2.06)}]
  \draw[black!30,fill=black!3,rounded corners=3pt] (0,0) rectangle (3.95,1.98);
  \node[sgfont=6.2,text=black!55] at (1.7,1.72) {조건 = 색 · 모양 · 수};
  \node[sgfont=6.2,text=sgorange!85!black] at (3.16,1.72) {액션};
  \foreach \k/\y in {if/1.36,elif/1.0,elif/.64,else/.28}{%
    \node[sgfont=7,anchor=west,font=\ttfamily\fontsize{7}{8}\selectfont,text=sgblue!80!black] at (.18,\y) {\k};
    \ifthenelse{\equal{\k}{else}}{}{\fill[black!22,rounded corners=1pt] (1.05,\y-.11) rectangle (2.35,\y+.11);}
    \fill[sgorange!35,rounded corners=1pt] (2.7,\y-.11) rectangle (3.62,\y+.11);}
\end{scope}
\draw[black!35,line width=.6pt,-{Stealth[length=3.5pt]}] (3.7,3.05)--(4.18,3.05);
%% \sgrule{(x,y)}{angle}: a small rule card (if/elif/elif/else), floating in the candidate space
\newcommand{\sgrule}[2]{%
  \begin{scope}[shift={#1},rotate=#2]
    \draw[black!35,fill=white,rounded corners=1pt,line width=.4pt] (-.2,-.17) rectangle (.2,.17);
    \foreach \y in {.093,.031,-.031}{\fill[black!25] (-.15,\y-.018) rectangle (.02,\y+.018);
      \fill[sgorange!45] (.06,\y-.018) rectangle (.15,\y+.018);}
    \fill[sgorange!45] (.06,-.111) rectangle (.15,-.075);
  \end{scope}}
%% ---- bottom: the new signal, top centre
\draw[black!30,dashed,line width=.4pt] (0,1.86)--(8.2,1.86);
\node[sgfont=7,text=black!70] at (4.1,1.64) {새 시그널 (정답 미공개)};
\draw[black!45,fill=white,rounded corners=2pt,line width=.8pt] (3.4,.98) rectangle (4.8,1.42);
\sgsig{(4.1,1.2)}{sgblue}{square}{3}
%% ---- bottom left: the solver infers its BELIEF p
\sgrobot[.62]{(1.57,1.16)}{sgorange}{solver}
\draw[black!40,line width=.6pt,-{Stealth[length=3.5pt]}] (3.32,1.2)--(2.02,1.2);
\node[sgfont=6.3,text=black!60,anchor=north] at (2.67,1.17) {추론};
\draw[sgorange!80,line width=.6pt,-{Stealth[length=3.5pt]}] (1.57,.86)--(1.57,.52);
%% ---- bottom right: rule candidates, filtered by the twenty clues
\node[sgfont=7,text=black!70,anchor=east] at (8.15,1.64) {숨은 규칙 후보};
\fill[black!4] (6.8,1.12) ellipse (1.3 and .36);
\draw[black!22,line width=.4pt,dash pattern=on 1.5pt off 1pt] (6.8,1.12) ellipse (1.3 and .36);
\foreach \x/\y/\r/\ok in {5.88/1.14/-8/1, 6.34/.98/6/0, 6.8/1.2/-4/1, 7.26/.99/8/1, 7.72/1.15/-6/0}{%
  \sgrule{(\x,\y)}{\r}
  \ifthenelse{\equal{\ok}{1}}{\sgcheck[.55]{(\x+.2,\y+.17)}}{\sgcross[.55]{(\x+.2,\y+.17)}}}
\draw[black!40,line width=.6pt,-{Stealth[length=3.5pt]}] (4.88,1.2)--(5.47,1.1);
\draw[black!60,line width=.6pt,-{Stealth[length=3.5pt]}] (6.62,.75)--(6.62,.52);
%% ---- bottom: the two distributions side by side (bars over the action glyphs), compared
\foreach \x/\sty/\a/\b/\c/\d in {.45/o/.10/.50/.10/.30, 5.5/w/0/.70/0/.30}{%
  \foreach \v/\act [count=\i from 0] in {\a/L, \b/R, \c/S, \d/J}{%
    \pgfmathsetmacro\bx{\x+\i*.62}%
    \ifthenelse{\equal{\sty}{o}}{\fill[sgorange!75] (\bx,-.42) rectangle ++(.38,\v);}%
      {\draw[black!75,line width=.7pt] (\bx,-.42) rectangle ++(.38,\v);}
    \sgact{(\bx+.19,-.56)}{\act}}
  \draw[black!55,line width=.5pt] (\x-.1,-.42)--(\x+2.34,-.42);}
\node[sgfont=6.6] at (1.57,.4) {solver의 BELIEF $p$};
\node[sgfont=6.6] (sgstitle) at (6.7,.4) {규칙들이 고른 비율 $s$};
\sgcheck[.5]{($(sgstitle.west)+(-.08,.01)$)}
%% the comparison in the middle
\draw[black!80,line width=.8pt,{Stealth[length=4pt]}-{Stealth[length=4pt]}] (2.9,-.06)--(5.3,-.06);
\node[sgfont=6.6,text=black!75] at (4.1,.12) {두 분포의 거리};
\node[font=\sffamily\bfseries\fontsize{7}{8}\selectfont,text=black!90] at (4.1,-.26) {= 추론 능력};
\end{tikzpicture}
\endgroup
}
\caption{시그널 풀이 과정. $s$는 단서에 맞는 규칙의 액션 비율, $p$는 기록된 BELIEF를 정규화한 값이다.}
\label{fig:signal}
\end{figure}
